% =============================================================
% บทที่ 4 การใช้งานสำหรับนักท่องเที่ยว (ผู้ใช้งานทั่วไปและสมาชิก)
% เนื้อหาเรียงตาม Use Case ในบทที่ 4 ของเล่มรายงาน
% ตรวจกับเว็บจริงที่รันบน localhost และซอร์สโค้ดใน main repo/frontend (27 ต.ค. 2569)
% =============================================================
\label{ch:tourist}

\section{การค้นหาสถานที่ท่องเที่ยว}
\label{sec:search-places}

หน้ารายการสถานที่ท่องเที่ยวแสดงดังภาพที่~\ref{fig:places-list}

\appscreen{places-list}{หน้ารายการสถานที่ท่องเที่ยว}

\begin{steps}
  \item คลิกเมนู \btn{สถานที่ท่องเที่ยว} หรือพิมพ์คำค้นในช่องค้นหาของหน้าแรก
  \item พิมพ์ชื่อสถานที่ใน \field{ช่องค้นหา} หรือเลือกหมวดหมู่ ได้แก่ ทั้งหมด คาเฟ่ วัด พิพิธภัณฑ์ สวนสาธารณะ อุทยานแห่งชาติ ตลาด ห้างสรรพสินค้า ร้านของฝาก สถานที่ท่องเที่ยว และร้านอาหาร
  \item เลือกการเรียงลำดับตามความแนะนำ คะแนนสูงสุด หรือยอดนิยม
        \result{ระบบแสดงการ์ดสถานที่ที่ตรงกับเงื่อนไข ถ้าไม่พบ ระบบแนะนำให้ลองค้นด้วยคำอื่น}
\end{steps}

\section{การดูรายละเอียดสถานที่}

คลิกการ์ดสถานที่เพื่อเปิดหน้ารายละเอียดดังภาพที่~\ref{fig:place-detail} ซึ่งแสดงหมวดหมู่ คะแนนและช่วงราคา แกลเลอรีรูปภาพ สรุปโดย Dino AI สิ่งอำนวยความสะดวก เช่น ที่จอดรถ การเข้าถึงสำหรับผู้พิการ และสถานีชาร์จรถยนต์ไฟฟ้า เวลาทำการ ที่อยู่และแผนที่ ข้อมูลติดต่อ และรีวิวจากผู้เยี่ยมชม ท้ายหน้าแสดงวันที่เพิ่มข้อมูลและวันที่อัปเดตล่าสุด

\appscreen{place-detail}{หน้ารายละเอียดสถานที่ท่องเที่ยว}

ปุ่มในหน้ารายละเอียด ได้แก่ \btn{เปิดเส้นทาง Google Maps} \btn{โทรติดต่อ} \btn{เขียนรีวิวบน Google} และ \btn{แชร์} ซึ่งคัดลอกลิงก์ของหน้านี้ สถานที่ที่ปิดชั่วคราวหรือปิดถาวรจะมีป้ายสถานะกำกับ

\begin{note}
ระบบไม่มีการเขียนรีวิวภายในเว็บไซต์ รีวิวที่แสดงมาจาก Google Maps ถ้าข้อมูลไม่ตรงกับความเป็นจริง สมาชิกแจ้งได้ตามหัวข้อ~\ref{sec:report-data}
\end{note}

\section{การดูกิจกรรมและเทศกาล}
\label{sec:events}

เมนู \btn{กิจกรรม \& เทศกาล} แสดงรายการกิจกรรมพร้อมช่องค้นหาและตัวกรองสถานะ ดังภาพที่~\ref{fig:events-list} ระบบกำหนดสถานะจากวันที่จัดงานให้อัตโนมัติตามตารางที่~\ref{tab:event-status}

\appscreen{events-list}{หน้ารายการกิจกรรมและเทศกาล}

\begin{table}[H]
\caption{สถานะของกิจกรรม}
\label{tab:event-status}
\centering
\begin{tabular}{|L{3.2cm}|L{10cm}|}
\hline
\textbf{สถานะ} & \textbf{ความหมาย} \\ \hline
เร็ว ๆ นี้ & ยังไม่ถึงวันจัดงาน \\ \hline
กำลังจัดอยู่ & อยู่ในช่วงวันจัดงาน \\ \hline
จบแล้ว & เลยวันจัดงานแล้ว แชตบอตจะไม่นำไปแนะนำ \\ \hline
ยกเลิก & ผู้จัดงานยกเลิกงาน ยังเปิดดูรายละเอียดได้ \\ \hline
\end{tabular}
\end{table}

คลิกกิจกรรมเพื่อดูรายละเอียดดังภาพที่~\ref{fig:event-detail} ได้แก่ ภาพปกงาน คำอธิบาย วันจัดงาน สถานที่จัดงานพร้อมแผนที่ ค่าเข้าชม ผู้จัดงาน และกลุ่มที่เหมาะสม

\appscreen{event-detail}{หน้ารายละเอียดกิจกรรม}

\section{การสอบถามข้อมูลกับแชตบอต}
\label{sec:chatbot}

หน้าต่างแชตบอตแสดงดังภาพที่~\ref{fig:chatbot}

\appscreen[0.9]{chatbot}{หน้าต่างแชตบอตน้องไดโน}

\begin{steps}
  \item คลิกปุ่มรูปน้องไดโนที่มุมขวาล่างของหน้าใดก็ได้
  \item พิมพ์คำถามเกี่ยวกับสถานที่ อาหาร การเดินทาง หรือเทศกาลในช่อง \field{พิมพ์คำถามของคุณ...} แล้วคลิกปุ่มส่ง
        \result{แชตบอตแสดงคำตอบต่อเนื่องทีละส่วน พร้อมการ์ดสถานที่ที่ใช้อ้างอิงใต้คำตอบ}
\end{steps}

ถ้าคำถามระบุพื้นที่ เช่น ``แถวท่าพระ'' แต่ในระบบมีสถานที่ในพื้นที่นั้นน้อย ระบบแสดงสถานะ ``กำลังค้นเพิ่มจาก Google...'' แล้วนำสถานที่จาก Google Maps มาตอบเพิ่ม การ์ดเหล่านั้นติดป้าย ``แหล่งจาก Google'' และมีปุ่ม \btn{แผนที่} \btn{โทร} และ \btn{ค้นหา} ดังภาพที่~\ref{fig:chatbot-google}

\appscreen[0.4]{chatbot-google}{คำตอบของแชตบอตที่ค้นเพิ่มจาก Google}

\begin{note}
ข้อมูลที่มาจาก Google ยังไม่ผ่านการตรวจสอบของทีมงาน คำตอบจึงมีหมายเหตุกำกับไว้ ถ้าไม่พบข้อมูลที่เกี่ยวข้อง แชตบอตจะแจ้งว่าไม่มีข้อมูลแทนการเดาคำตอบ
\end{note}

\section{การวางแผนทริปด้วย AI}
\label{sec:trip-plan}

แบบฟอร์มวางแผนทริปมี 4 ขั้นตอน ได้แก่ การเดินทาง ความสนใจ งบประมาณ และช่วงเวลา ขั้นตอนแรกแสดงดังภาพที่~\ref{fig:trip-step1}

\begin{screenshot}{images/screens/trip-step1.jpg}{แบบฟอร์มวางแผนทริป ขั้นตอนที่ 1 การเดินทาง}{fig:trip-step1}
  \callout{0.265}{0.225}{1}
  \callout{0.265}{0.355}{2}
  \callout{0.265}{0.545}{3}
  \callout{0.72}{0.48}{4}
  \callout{0.265}{0.955}{5}
\end{screenshot}

\begin{calloutlist}
  \item \field{เริ่มเที่ยววันไหน} เลือก วันนี้ พรุ่งนี้ หรือเสาร์นี้ หรือเลือกวันที่เอง
  \item \field{เที่ยวกี่วัน} 1 ถึง 5 วัน ระบบแสดงช่วงวันที่ของทริปให้ตรวจสอบ
  \item \field{แต่ละวันออกเดินทางจากที่ไหน} พิมพ์ค้นหาที่พักหรือปักหมุดบนแผนที่ เว้นว่างได้ถ้ายังไม่ได้จองที่พัก
  \item ปุ่ม \btn{ตำแหน่งปัจจุบัน} ใช้ตำแหน่งของอุปกรณ์เป็นจุดออกเดินทาง
  \item \field{วันสุดท้ายจบที่ไหน} กลับจุดเริ่มต้น สนามบินขอนแก่น บขส. แห่งที่ 3 สถานีรถไฟขอนแก่น อำเภออื่น หรือปักหมุดเอง
\end{calloutlist}

\begin{steps}
  \item คลิกเมนู \btn{วางแผนทริป AI}
  \item ขั้นตอนที่ 1 กรอกข้อมูลการเดินทางตามภาพที่~\ref{fig:trip-step1} แล้วคลิกปุ่ม \btn{ถัดไป}
        \begin{enumerate}
          \item ถ้าเว้นจุดออกเดินทางไว้ ระบบเริ่มแต่ละวันจากใจกลางเมืองขอนแก่น และแนะนำย่านที่ควรพักให้หลังสร้างแผน
          \item ถ้าวันสุดท้ายจบที่สนามบิน สถานีขนส่ง หรือสถานีรถไฟ ให้ระบุเวลาที่เครื่องหรือรถออก ระบบจะจัดวันสุดท้ายให้ไปถึงทันเวลา
        \end{enumerate}
  \item ขั้นตอนที่ 2 เลือกความสนใจได้หลายข้อ (ภาพที่~\ref{fig:trip-step2}) ได้แก่ ธรรมชาติ วัฒนธรรม/ศาสนา ไดโนเสาร์ คาเฟ่ ช้อปปิ้ง/หัตถกรรม และครอบครัว เลือกประเภทอาหารใน \field{อยากกินอะไร?} (ไม่บังคับ) และพิมพ์สถานที่ที่ต้องการไปแน่นอนใน \field{มีสถานที่ที่อยากไปแน่ๆ มั้ย?} (ไม่บังคับ)
  \item ขั้นตอนที่ 3 เลือกงบประมาณ (ประหยัด ปานกลาง หรือหรูหรา) จังหวะการเที่ยว (สายชิลล์ กำลังดี หรือสายลุย) และขอบเขตพื้นที่ (ภาพที่~\ref{fig:trip-step3}) ขอบเขต ``เมือง'' ใช้เฉพาะอำเภอเมืองขอนแก่น ส่วน ``ทั่วขอนแก่น'' รวมสถานที่รอบนอก เช่น ภูเวียง อุบลรัตน์
  \item ขั้นตอนที่ 4 กำหนดเวลาเริ่มและสิ้นสุดของแต่ละวัน (ค่าเริ่มต้น 09:00 ถึง 18:00 น.) แล้วคลิกปุ่ม \btn{สร้างแผนการเดินทาง} (ภาพที่~\ref{fig:trip-step4})
        \result{ระบบแสดงหน้ากำลังประมวลผล (ภาพที่~\ref{fig:trip-loading}) แล้วแสดงแผนการเดินทางรายวัน (ภาพที่~\ref{fig:trip-result})}
\end{steps}

\appscreen{trip-step2}{แบบฟอร์มวางแผนทริป ขั้นตอนที่ 2 ความสนใจ}

\appscreen{trip-step3}{แบบฟอร์มวางแผนทริป ขั้นตอนที่ 3 งบประมาณ}

\appscreen{trip-step4}{แบบฟอร์มวางแผนทริป ขั้นตอนที่ 4 ช่วงเวลา}

\appscreen[0.7]{trip-loading}{หน้าสถานะกำลังประมวลผลแผนการเดินทาง}

\begin{note}
ประเภทอาหารใช้เลือกร้านมื้อกลางวันและมื้อเย็น ส่วน ``อาหารเช้า'' จะเพิ่มมื้อเช้าให้วันที่เริ่มเที่ยวก่อน 09:45 น. ทริปหลายวันอาจใช้เวลาประมวลผลเกิน 1 นาที
\end{note}

หน้าผลแผนการเดินทางแสดงดังภาพที่~\ref{fig:trip-result}

\begin{screenshot}{images/screens/trip-result.jpg}{หน้าแสดงผลแผนการเดินทาง}{fig:trip-result}
  \callout{0.105}{0.30}{1}
  \callout{0.105}{0.755}{2}
  \callout{0.105}{0.87}{3}
  \calloutline{0.40}{0.115}{0.40}{0.19}{4}
  \calloutline{0.79}{0.115}{0.79}{0.19}{5}
  \callout{0.89}{0.45}{6}
  \callout{0.89}{0.80}{7}
\end{screenshot}

\begin{calloutlist}
  \item สรุปทริป ได้แก่ จำนวนวัน ระยะทางรวม และพอยท์ที่จะได้จากสถานที่ในแผน
  \item เงื่อนไขที่เลือกไว้ในแบบฟอร์ม
  \item ``ทำไมน้องไดโนจัดแบบนี้'' อธิบายเหตุผลของแผน
  \item วันที่ของแผน (ทริปหลายวันแสดงทีละวัน)
  \item ปุ่ม \btn{เปิดใน Google Maps} เปิดเส้นทางของวันนั้น
  \item แผนที่เส้นทางพร้อมลำดับจุดแวะ
  \item ตารางเวลา ได้แก่ เวลาที่ถึงแต่ละสถานที่ ระยะทางและเวลาเดินทางระหว่างจุด และค่าใช้จ่ายโดยประมาณ
\end{calloutlist}

ถ้าต้องการเปลี่ยนเงื่อนไข ให้คลิกปุ่ม \btn{สร้างแผนใหม่} เพื่อกลับไปยังแบบฟอร์ม ส่วน ``ข้อควรทราบเกี่ยวกับแผนนี้'' แจ้งข้อจำกัดของแผน เช่น สถานที่ที่อาจปิดในวันนั้น

ถ้าทริปมีตั้งแต่ 2 วันขึ้นไปและไม่ได้ระบุจุดออกเดินทาง ระบบคำนวณย่านที่ควรเข้าพักจากสถานที่ในแผน แล้วแสดงกล่อง ``ยังไม่มีที่พัก? พักแถวนี้สะดวกที่สุด'' ดังภาพที่~\ref{fig:trip-stay} พร้อมระยะห่างโดยประมาณจากจุดเริ่มและจุดจบของแต่ละวัน คลิกปุ่ม \btn{ดูที่พักแถวนี้ใน Google Maps} เพื่อค้นหาที่พักในย่านนั้น

\appscreen[0.45]{trip-stay}{กล่องแนะนำย่านที่พักในหน้าผลแผนการเดินทาง}

เมื่อจองที่พักแล้ว ให้คลิกปุ่ม \btn{สร้างแผนใหม่} แล้วปักหมุดที่พักเป็นจุดออกเดินทาง เวลาในแผนจะแม่นยำขึ้น

\begin{warning}
ผู้ใช้งานทั่วไปที่ไม่ได้เข้าสู่ระบบ แผนจะหายเมื่อปิดหน้านี้ ควรเข้าสู่ระบบก่อนสร้างแผน ระบบจะบันทึกแผนไว้ในเมนู \btn{ทริปของฉัน} ให้อัตโนมัติ
\end{warning}

\section{การจัดการทริปของฉัน}
\label{sec:my-trips}

หัวข้อนี้ต้องเข้าสู่ระบบก่อน แผนที่สร้างขณะเข้าสู่ระบบจะบันทึกไว้ในเมนู \btn{ทริปของฉัน} เรียงจากแผนล่าสุด ดังภาพที่~\ref{fig:my-trips} ค้นหาตามชื่อทริปหรือกรองเฉพาะรายการโปรดได้

\appscreen{my-trips}{หน้าทริปของฉัน}

การจัดการแต่ละแผนทำได้ดังนี้

\begin{enumerate}
\item คลิกแผนเพื่อเปิดหน้าผลแผนการเดินทางอีกครั้ง โดยระบบไม่ต้องสร้างแผนใหม่
\item ปุ่ม \btn{เปลี่ยนชื่อ} ตั้งชื่อแผนได้ไม่เกิน 120 ตัวอักษร
\item คลิกรูปดาวเพื่อเพิ่มหรือเอาออกจากรายการโปรด
\item ปุ่ม \btn{ทำสำเนา} สร้างแผนใหม่ที่มีเนื้อหาเดียวกัน โดยต่อท้ายชื่อว่า ``(สำเนา)''
\item ปุ่ม \btn{ลบ} ลบแผนพร้อมรายการวันและจุดแวะทั้งหมด ระบบขอให้ยืนยันก่อนลบ
\item ปุ่ม \btn{สร้างแผนใหม่} ไปยังแบบฟอร์มวางแผนทริป
\end{enumerate}

\begin{important}
แต่ละบัญชีบันทึกแผนได้ไม่เกิน 50 แผน เมื่อครบแล้วจะทำสำเนาไม่ได้ ต้องลบแผนเก่าก่อน
\end{important}

\section{การสะสมพอยท์ด้วย QR Code}
\label{sec:points}

เมนู \btn{พอยท์สะสม} ของผู้ใช้งานทั่วไปแสดงบัตรสะสมพอยท์ สถานที่ร่วมรายการพร้อมจำนวนพอยท์ ขั้นตอนการสะสม และของรางวัลที่แลกได้ ดังภาพที่~\ref{fig:points-guest} การสะสมพอยท์ต้องเข้าสู่ระบบก่อน โดยคลิกปุ่ม \btn{สมัครสมาชิกฟรี} หรือ \btn{เข้าสู่ระบบ} ในหน้านี้

\appscreen{points-guest}{หน้าพอยท์สะสมของผู้ใช้งานทั่วไป}

เมื่อเข้าสู่ระบบแล้ว หน้า ``พอยท์สะสมของคุณ'' แสดงพอยท์คงเหลือ ปุ่มสแกน QR Code สถานที่ที่มี QR รับพอยท์ และของรางวัลพร้อมจำนวนพอยท์ที่ใช้ ของรางวัลที่หมดสต็อกมีป้าย ``ของหมด'' ดังภาพที่~\ref{fig:points-member}

\appscreen{points-member}{หน้าพอยท์สะสมของสมาชิก}

\begin{steps}
  \item ไปยังสถานที่ร่วมรายการ แล้วคลิกปุ่มสแกนในหน้าพอยท์สะสม หรือใช้กล้องของโทรศัพท์มือถือสแกนป้าย QR Code โดยตรง
  \item อนุญาตให้เบราว์เซอร์ใช้กล้องและตำแหน่งที่ตั้ง (ภาพที่~\ref{fig:qr-scan})
  \item เล็งกล้องไปที่ป้าย QR Code
        \result{ระบบเพิ่มพอยท์เข้าบัญชีทันทีและแสดงจำนวนพอยท์ที่ได้รับ}
\end{steps}

\appscreen[0.7]{qr-scan}{หน้าต่างสแกน QR Code}

\begin{important}
ระบบตรวจว่าอุปกรณ์อยู่ในรัศมีที่กำหนดของสถานที่จึงจะให้พอยท์ และ QR Code แต่ละจุดรับพอยท์ได้ครั้งเดียวต่อบัญชี
\end{important}

ถ้าสแกนไม่ผ่าน ระบบแสดง ``สแกนไม่สำเร็จ'' พร้อมเหตุผลตามตารางที่~\ref{tab:qr-errors}

\begin{table}[H]
\caption{ข้อความเมื่อสแกน QR Code ไม่สำเร็จ}
\label{tab:qr-errors}
\centering
\begin{tabular}{|L{5.4cm}|L{3.8cm}|L{4cm}|}
\hline
\textbf{ข้อความ} & \textbf{สาเหตุ} & \textbf{สิ่งที่ควรทำ} \\ \hline
คุณเคยสแกน QR นี้ไปแล้ว & บัญชีเคยรับพอยท์จาก QR นี้แล้ว & สะสมพอยท์จากสถานที่อื่น \\ \hline
ไม่พบ QR นี้ / QR นี้ปิดใช้งานอยู่ / QR นี้หมดอายุแล้ว & QR ใช้งานไม่ได้ & แจ้งเจ้าหน้าที่ของสถานที่ \\ \hline
ต้องเปิดการเข้าถึงตำแหน่งเพื่อสแกน QR นี้ & ไม่ได้อนุญาตตำแหน่งที่ตั้ง & อนุญาตตำแหน่งในการตั้งค่าเบราว์เซอร์แล้วสแกนใหม่ \\ \hline
ข้อความแจ้งระยะห่างจากสถานที่ & อุปกรณ์อยู่นอกรัศมีของสถานที่ & เข้าใกล้สถานที่แล้วสแกนใหม่ \\ \hline
บัญชีนี้ไม่สามารถใช้งานได้ กรุณาติดต่อผู้ดูแลระบบ & บัญชีถูกระงับ & ติดต่อผู้ดูแลระบบ \\ \hline
\end{tabular}
\end{table}

\section{การแลกของรางวัล}

หัวข้อนี้ต้องเข้าสู่ระบบก่อน การแลกของรางวัลทำที่เคาน์เตอร์ของสำนักงาน ททท. ขอนแก่น

\begin{steps}
  \item แจ้งชื่อหรือเบอร์โทรศัพท์ที่ใช้สมัครกับเจ้าหน้าที่ พร้อมของรางวัลที่ต้องการ
        \result{เจ้าหน้าที่ตรวจพอยท์ หักพอยท์ และมอบของรางวัลผ่านหน้าผู้ดูแลระบบ (หัวข้อ~\ref{sec:admin-redeem})}
  \item ตรวจรายการแลกได้ที่ \menu{โปรไฟล์ของฉัน > ประวัติ}
\end{steps}

\section{การแจ้งข้อมูลไม่ถูกต้อง}
\label{sec:report-data}

หัวข้อนี้ต้องเข้าสู่ระบบก่อน ผู้ใช้งานทั่วไปที่คลิกปุ่มนี้จะได้รับแจ้งให้เข้าสู่ระบบ หน้าต่างแจ้งข้อมูลแสดงดังภาพที่~\ref{fig:report-data}

\appscreen[0.7]{report-data}{หน้าต่างแจ้งข้อมูลไม่ถูกต้อง}

\begin{steps}
  \item ในหน้ารายละเอียดสถานที่หรือกิจกรรม คลิกปุ่ม \btn{แจ้งข้อมูลไม่ถูกต้อง}
  \item เลือกหัวข้อที่ไม่ถูกต้องอย่างน้อย 1 ข้อ
        \begin{enumerate}
          \item สถานที่: ชื่อ ที่อยู่ ตำแหน่ง เวลาทำการ เบอร์โทรศัพท์ เว็บไซต์ ปิดกิจการ หรือรูปภาพ
          \item กิจกรรม: ชื่ออีเวนต์ วันและเวลาจัดงาน สถานที่จัดงาน ค่าเข้าชม สถานะยกเลิกหรือเลื่อน ผู้จัดงาน หรือรูปภาพ
        \end{enumerate}
  \item พิมพ์ \field{รายละเอียดเพิ่มเติม} ไม่เกิน 500 ตัวอักษร แล้วคลิกส่ง
        \result{ระบบแสดง ``ส่งรายงานแล้ว ขอบคุณที่ช่วยกันปรับปรุงข้อมูล'' และรายงานเข้าคิวของผู้ดูแลระบบ}
\end{steps}

\begin{note}
ถ้าเลือกหัวข้อ ``อื่นๆ'' ต้องอธิบายปัญหาในช่องรายละเอียดเพิ่มเติมด้วย หัวข้อที่เคยแจ้งแล้วและยังรอตรวจสอบจะแจ้งซ้ำไม่ได้
\end{note}

\section{การจัดการโปรไฟล์}
\label{sec:profile}

หัวข้อนี้ต้องเข้าสู่ระบบก่อน คลิก \menu{เมนูผู้ใช้ > โปรไฟล์ของฉัน} จะพบหน้าโปรไฟล์ซึ่งมี 5 แท็บ ดังภาพที่~\ref{fig:profile-overview}

\appscreen{profile-overview}{หน้าโปรไฟล์ แท็บภาพรวม}

\begin{enumerate}
\item แท็บภาพรวม ข้อมูลสรุปของผู้ใช้ พอยท์คงเหลือ และปุ่มลัด \btn{สแกน QR} และ \btn{ดูประวัติพอยท์} คลิกปุ่ม \btn{เปลี่ยนรูป} เพื่ออัปโหลดรูปโปรไฟล์ (jpg png webp หรือ gif ไม่เกิน 2 MB) แล้วลากรูปเพื่อเลือกส่วนที่แสดงในวงกลม
\item แท็บข้อมูล แก้ไขข้อมูลส่วนตัว แล้วคลิกปุ่ม \btn{บันทึกการเปลี่ยนแปลง} (ภาพที่~\ref{fig:profile-info})
\item แท็บประวัติ ประวัติพอยท์จากการสแกน QR การแลกรางวัล และการปรับพอยท์โดยผู้ดูแลระบบ (ภาพที่~\ref{fig:profile-history}) รายการแลกที่ถูกยกเลิกแสดงว่า ``ยกเลิกแล้ว (คืนพอยท์)''
\item แท็บกิจกรรมของฉัน ส่งคำขอเพิ่มกิจกรรมและติดตามสถานะ (หัวข้อ~\ref{sec:event-request})
\item แท็บความปลอดภัย เปลี่ยนรหัสผ่าน เปลี่ยนอีเมล และลบบัญชี (ภาพที่~\ref{fig:profile-security})
\end{enumerate}

\appscreen{profile-info}{หน้าโปรไฟล์ แท็บข้อมูล}

\appscreen{profile-history}{หน้าโปรไฟล์ แท็บประวัติ}

\appscreen{profile-security}{หน้าโปรไฟล์ แท็บความปลอดภัย}

การเปลี่ยนรหัสผ่านต้องกรอกรหัสผ่านปัจจุบัน เมื่อเปลี่ยนแล้วอุปกรณ์อื่นที่เข้าสู่ระบบอยู่จะถูกออกจากระบบ การเปลี่ยนอีเมลต้องกรอกอีเมลใหม่และรหัสผ่าน ระบบส่งลิงก์ยืนยันไปยังอีเมลใหม่ และเปลี่ยนอีเมลเมื่อเปิดลิงก์ภายใน 24 ชั่วโมง

\begin{warning}
การลบบัญชีต้องกรอกรหัสผ่านเพื่อยืนยัน บัญชีจะเข้าสู่ระบบไม่ได้อีก พอยท์ที่เหลือนำไปแลกรางวัลไม่ได้ และอีเมลนี้นำไปสมัครใหม่ไม่ได้ การกู้คืนบัญชีต้องติดต่อผู้ดูแลระบบ
\end{warning}

\section{การส่งคำขอเพิ่มกิจกรรม}
\label{sec:event-request}

หัวข้อนี้ต้องเข้าสู่ระบบก่อน สมาชิกเสนอกิจกรรมที่ยังไม่มีในระบบได้จากแท็บกิจกรรมของฉัน ดังภาพที่~\ref{fig:profile-events}

\appscreen{profile-events}{หน้าโปรไฟล์ แท็บกิจกรรมของฉัน}

\begin{steps}
  \item คลิก \menu{โปรไฟล์ของฉัน > กิจกรรมของฉัน} แล้วคลิกปุ่ม \btn{ส่งกิจกรรมใหม่}
  \item กรอก \field{ชื่องาน} \field{ประเภทงาน} \field{ช่วงวันจัดงาน} \field{ชื่อสถานที่หรือสนาม} (เชื่อมกับสถานที่ในระบบได้) \field{ค่าเข้างาน} \field{ผู้จัดงาน} \field{กลุ่มที่เหมาะสม} และ \field{รายละเอียดงาน} พร้อมแนบภาพงาน
  \item คลิกส่ง
        \result{คำขอมีสถานะรออนุมัติ และแสดงบนหน้าเว็บเมื่อผู้ดูแลระบบอนุมัติแล้วเท่านั้น}
\end{steps}

\begin{tip}
คัดลอกข้อความจากโพสต์เพจ Facebook ของงานมาวาง แล้วคลิกปุ่ม \btn{ดึงข้อมูลอัตโนมัติ} ให้ AI ช่วยเติมแบบฟอร์ม และตรวจข้อมูลก่อนส่งทุกครั้ง
\end{tip}

สถานะของคำขอแสดงดังตารางที่~\ref{tab:event-request-status} แต่ละบัญชีมีคำขอที่รออนุมัติได้ไม่เกิน 5 รายการ

\begin{table}[H]
\caption{สถานะของคำขอเพิ่มกิจกรรม}
\label{tab:event-request-status}
\centering
\begin{tabular}{|L{3.2cm}|L{4.6cm}|L{5.4cm}|}
\hline
\textbf{สถานะ} & \textbf{ความหมาย} & \textbf{สิ่งที่สมาชิกทำได้} \\ \hline
รออนุมัติ & รอผู้ดูแลระบบตรวจสอบ & แก้ไขหรือยกเลิกคำขอ \\ \hline
อนุมัติแล้ว & ระบบสร้างกิจกรรมบนหน้าเว็บแล้ว & ดูกิจกรรมในเมนูกิจกรรม \& เทศกาล \\ \hline
ไม่อนุมัติ & ผู้ดูแลระบบปฏิเสธคำขอ & อ่านเหตุผลที่ผู้ดูแลระบบระบุ แล้วส่งคำขอใหม่ถ้าแก้ไขได้ \\ \hline
\end{tabular}
\end{table}

\section{การตอบแบบสอบถามความพึงพอใจ}
\label{sec:survey}

ผู้ใช้ทุกคนตอบแบบสอบถามได้โดยไม่ต้องเข้าสู่ระบบ แบบสอบถามแสดงดังภาพที่~\ref{fig:survey}

\appscreen{survey}{หน้าแบบสอบถามความพึงพอใจ}

\begin{steps}
  \item คลิกปุ่ม \btn{ประเมินความพึงพอใจ} ที่ส่วนท้ายของเว็บไซต์ หรือ \btn{ทำแบบประเมิน} ในหน้าผลแผนการเดินทาง
  \item ตอบแบบสอบถาม 4 ตอน ได้แก่ ข้อมูลทั่วไป ประสบการณ์ก่อนใช้งาน Dino ความพึงพอใจต่อระบบ (คะแนน 5 ระดับ 16 ข้อ) และข้อเสนอแนะ (ไม่บังคับ)
  \item คลิกปุ่ม \btn{ส่งแบบสอบถาม}
        \result{ระบบบันทึกคำตอบแบบไม่ระบุตัวตนและแสดง ``ขอบคุณที่ร่วมตอบแบบสอบถาม''}
\end{steps}

\begin{note}
คำถามเกี่ยวกับแชตบอตและการวางแผนเที่ยวเลือกตอบ ``ยังไม่เคยใช้'' ได้ ระบบไม่นำข้อนั้นไปคำนวณคะแนน แบบสอบถามตอบได้เครื่องละ 1 ครั้ง
\end{note}
